\color{ink}\frontmatter
\begin{titlepage}\centering
\vspace*{2.1cm}
{\Huge\sffamily\bfseries\color{navy} MAT0339\par}\vspace{0.3cm}
{\LARGE\sffamily\color{ink} Cahier d'exercices et de problèmes\par}\vspace{0.65cm}
{\large\sffamily\color{teal} Un parcours aligné sur chaque section du manuel\par}\vspace{1.3cm}
\begin{tcolorbox}[width=0.76\textwidth,colback=mistteal,borderline west={3pt}{0pt}{teal}]
\centering\sffamily
Comprendre \textbullet\ Calculer \textbullet\ Représenter \textbullet\ Interpréter \textbullet\ Vérifier
\end{tcolorbox}
\vfill
{\Large\sffamily\bfseries\color{navy} Xia Xiao\par}\vspace{0.2cm}
{\large\sffamily\color{slate} Département de mathématiques, UQAM\par}\vspace{0.7cm}
{\large\sffamily\color{teal} Édition révisée -- 5 septembre 2026\par}
\end{titlepage}
\chapter*{Mode d'emploi}\addcontentsline{toc}{chapter}{Mode d'emploi}
Ce cahier accompagne le manuel \emph{MAT0339 - Mathématiques}. Il conserve les cinq parties, les dix-sept chapitres et les intitulés de sections du manuel. La révision de septembre précise quelques consignes du chapitre 3 sans renuméroter les exercices. Chaque section reçoit un exercice en trois volets :
\begin{enumerate}
\item une question de compréhension ou de représentation;
\item un travail technique ciblé;
\item une application, une interprétation ou un contrôle.
\end{enumerate}
\begin{designbox}
Les exercices ne sont pas conçus comme une accumulation de calculs. Lorsque plusieurs méthodes sont possibles, il faut en choisir une, annoncer les hypothèses, conserver les unités, puis vérifier la plausibilité du résultat. Les problèmes de synthèse demandent de relier plusieurs représentations ou plusieurs sections du chapitre.
\end{designbox}
\ifcorrige Cette version contient des réponses brèves à la fin du cahier.\else Cette version étudiante ne contient pas les réponses; une version avec réponses est disponible séparément.\fi
\tableofcontents\mainmatter
